
\section{Imported background, live/reserve boundary, and companion concordance}\label{app:ledger-concordance}

This appendix fixes the boundary between the internal proof burden carried in the core paper and the imported or deferred material moved to companion papers. It also records the live-versus-reserve split required by the Phase-I freeze.

\subsection*{Live proof surface versus reserve materials}

\begin{quote}\small
The live proof burden of the hardened suite is carried only by the four compiled roots --- \texttt{core.tex}, \texttt{companion1.tex}, \texttt{companion2.tex}, and \texttt{companion3.tex} --- together with the theorem-facing appendices compiled inside them. Support PDFs, archival source facsimiles not compiled in those roots, source-only routing appendices, build notes, and editorial memoranda are reserve materials only. They add no mathematics, carry no unique proof step, and are not part of the proof burden.
\end{quote}

The reserve category currently includes the support PDFs under \texttt{support/}, the source-only routing appendices \texttt{source/appendices/appendixP\_referee\_audit\_packet.tex} and \texttt{source/appendices/appendixQ\_executive\_proof\_sketch.tex}, and the editorial/build ledgers in the repository root. The synchronized Phase-I routing references are Appendix~\ref{app:packet-docket} and Appendix~\ref{app:phase1-freeze} of the core paper.

\subsection*{Imported background ledger}
\begin{description}[leftmargin=2.2em,style=nextline]
\item[KP/Dobrushin entrance-to-clustering package.] Used in Sections~\ref{sec:uv} and \ref{sec:route1}, and in the downstream clustering clause of Section~\ref{sec:assembly}. It supplies the standard clustering machinery once one-shot entrance has been established; it is not a hidden substitute for the Route~1 transport chain.
\item[Standard OS reconstruction / analytic continuation / uniqueness theorem family.] Used in Section~\ref{sec:endpoint}, especially \cref{thm:euclidean-os-dossier,thm:wightman-reconstruction,thm:faithful-wilson-universality}. It upgrades the constructive Euclidean packet to the Wightman/Haag--Kastler reconstruction and provides the uniqueness clause used in faithful-Wilson universality.
\end{description}

\subsection*{Phase-VIII adversarial-verification packet}
The proof-kernel root now carries a dedicated adversarial-verification pack in its Appendices~C--D, and the core paper now carries a synchronized Appendix~D recording the same internal sign-off scaffold at the public-surface level. These documents add no mathematics and do not alter the live theorem burden, but they freeze the four specialist attack tracks, the issue taxonomy, the kill criteria, and the final manuscript-level dispositions. The repository-root files \texttt{PHASE8\_TRACK1\_RG\_LATTICE\_REPORT.md}, \texttt{PHASE8\_TRACK2\_OS\_MEASURE\_REPORT.md}, \texttt{PHASE8\_TRACK3\_OPERATOR\_DOMAIN\_REPORT.md}, \texttt{PHASE8\_TRACK4\_ENDPOINT\_LOCALNET\_REPORT.md}, \texttt{PHASE8\_ISSUE\_TRACKER.md}, \texttt{PHASE8\_DISPOSITION\_MEMO.md}, and \texttt{PHASE8\_SIGNOFF\_CHECKLIST.md} synchronize the same issue IDs in lightweight editorial form.

\subsection*{Deferred modules}
\begin{itemize}
    \item Lane~B and the spectral-calculus wrapper remain downstream reformulations and do not occupy numbered theorem space in the core paper.
    \item Classical-family worked audits, sample-sector normalization checks, long constant ledgers, and archival Route~1 branches are moved entirely to companion papers or reserve support files.
    \item Kernel-parametrix detail, insertion-RG proofs, and long coefficient-matrix proofs are part of the full proof homes, not of the core theorem paper.
    \item AQFT commentary on sectors and field-algebra extensions remains downstream commentary and is not part of the theorem-level endpoint claim.
\end{itemize}

\subsection*{Proof-home concordance}
\begin{description}[leftmargin=2.2em,style=nextline]
\item[Core-only results.] New labels 1.2--1.5 and 2.1 remain in the core paper. The assembled main theorem 1.1 is proved in Core Section~\ref{sec:assembly}, with Companions~I--III cited theorem by theorem.
\item[Companion~I.] Core results 3.1, 3.2, 3.3, 4.1, 4.2, 4.3, 6.1, 6.2, 6.3, 6.4, 6.5, and 6.6 route to Companion~I, Sections~1--6.
\item[Companion~II.] Core results 5.1, 5.2, 5.3, 5.4, 5.5, 5.6, 5.7, 5.8, 5.9, 5.10, 5.11, and 5.12 route to Companion~II, Sections~1--6.
\item[Companion~III.] Core results 7.1, 7.2, 7.3, 7.4, 7.5, 7.6, 7.7, 7.8, 7.9, and 7.10 route to Companion~III, Sections~1--4.
\end{description}

This appendix is the editorial safety valve of the refactor: it keeps the body-facing theorem paper compact without obscuring where any proof, imported theorem family, or deferred module actually lives, and it now does so under an explicit live-versus-reserve boundary.
